% === Begin Label Data ===
% ID: 847
% 难度星级: 3
% 标签: 解三角形，三角恒等变换，余弦定理
% 备注: 
% 组卷引用次数: 12
% === End  Label Data ===

\begin{problem}{2026}{L}{}{}{解三角形}
已知 $\triangle ABC$ 的内角 $A, B, C$ 的对边分别为 $a, b, c$，且 $\sin(A+C)=8\sin^2\frac{B}{2}$.

（1）求 $\cos B$;

（2）若 $a+c=6$，$\triangle ABC$ 的面积为 2，求 $b$.
\end{problem}

\begin{answer}
$\displaystyle\dfrac{15}{17}; 2$
\end{answer}

\begin{solutions}
（1）在 $\triangle ABC$ 中，因为 $A+B+C=\pi$,所以 $\sin(A+C)=\sin(\pi-B)=\sin B$.

原式化为 $\sin B=8\sin^2\frac{B}{2}$.

利用二倍角公式得 $2\sin\frac{B}{2}\cos\frac{B}{2}=8\sin^2\frac{B}{2}$.

因为 $B\in(0,\pi)$,所以 $\sin\frac{B}{2}\neq 0$,两边同除以 $2\sin\frac{B}{2}$ 得 $\cos\frac{B}{2}=4\sin\frac{B}{2}$.

平方得 $\cos^2\frac{B}{2}=16\sin^2\frac{B}{2}$.

结合 $\sin^2\frac{B}{2}+\cos^2\frac{B}{2}=1$,解得 $\sin^2\frac{B}{2}=\displaystyle\frac{1}{17}$.

由降幂公式 $\cos B=1-2\sin^2\frac{B}{2}=1-\displaystyle\frac{2}{17}=\displaystyle\frac{15}{17}$.

（2）因为 $\cos B=\displaystyle\frac{15}{17}$ 且 $B\in(0,\pi)$,所以 $\sin B=\sqrt{1-\cos^2 B}=\displaystyle\frac{8}{17}$.

由面积公式 $S_{\triangle ABC}=\displaystyle\frac{1}{2}ac\sin B=2$,得 $ac=\displaystyle\frac{4}{\sin B}=\displaystyle\frac{17}{2}$.

由余弦定理 $b^2=a^2+c^2-2ac\cos B=(a+c)^2-2ac(1+\cos B)$.

代入 $a+c=6$,$ac=\displaystyle\frac{17}{2}$,$\cos B=\displaystyle\frac{15}{17}$ 得：

$$
b^2=6^2-2\times\displaystyle\frac{17}{2}\times\left(1+\displaystyle\frac{15}{17}\right)=36-32=4.
$$

所以 $b=\boxed{2}$.
\end{solutions}